\setcounter{section}{12}

\inputaufgabe
{}
{

Tunjukkan bahwa sebuah
\definitionsverweis {bundel vektor riil}{}{}
di atas satu titik
\zusatzklammer {yakni sebuah ruang topologis bertitik tunggal} {} {}
sama dengan sebuah ruang vektor riil berdimensi hingga.

}
{} {}

\inputaufgabe
{}
{

Misalkan
\maabb {p} { V } { X
} {}
sebuah
\definitionsverweis {bundel vektor riil}{}{}
di atas sebuah
\definitionsverweis {ruang topologis}{}{}
$X$. Tunjukkan bahwa $V$ merupakan sebuah
\definitionsverweis {ruang Hausdorff}{}{}
jika dan hanya jika $X$ merupakan ruang Hausdorff.

}
{} {}

\inputaufgabe
{}
{

Misalkan $X$ sebuah
\definitionsverweis {ruang topologis}{}{.}
Tunjukkan bahwa
\definitionsverweis {identitas}{}{}
\maabb {\operatorname{Id}_{  X  }} { X } { X
} {}
dapat dipandang sebagai sebuah
\definitionsverweis {bundel vektor riil}{}{}
ber-
\definitionsverweis {rank}{}{}
$0$.

}
{} {}

\inputaufgabe
{}
{

Tunjukkan bahwa bundel kernel dari
[[Reelle Ebene/Parameter/Kernbündel/Beispiel|Contoh 12.6]]
adalah trivial.

}
{} {}

\inputaufgabe
{}
{

Untuk
\definitionsverweis {bundel vektor}{}{}

\mavergleichskettealign
{\vergleichskettealign
{L
}
{ =} { \left\{ (r,s,t,u,v,w)  \mid  ru+sv+tw = 0 , \, (r,s,t) \neq (0,0,0)       \right\}
}
{ \subseteq} { (\R^3 \setminus \{(0,0,0)\}) \times \R^3
}
{ \longrightarrow} { \R^3 \setminus \{(0,0,0)\}
}
{ } {
}
}
{}{}{}
tentukan trivialisasi-trivialisasi linear di atas
\mathl{D(r)}{,}
\mathl{D(s)}{} dan
\mathl{D(t)}{,} yakni basis-basis yang bergantung pada $r,s,t$ di atas $D(r)$ dan seterusnya. Tentukan pemetaan-pemetaan perubahan basis pada

\mavergleichskette
{\vergleichskette
{D(rs)
}
{ = }{ D(r) \cap D(s)
}
{  }{
}
{    }{
}
{   }{
}
}
{}{}{.}

}
{} {}

\inputaufgabe
{}
{

Untuk
\definitionsverweis {bundel vektor}{}{}

\mavergleichskettealign
{\vergleichskettealign
{L
}
{ =} { \left\{ (r,s,t,u,v,w)  \mid  ru+sv+tw = 0 , \, (r,s,t) \neq (0,0,0)       \right\}
}
{ \subseteq} { (\R^3 \setminus \{(0,0,0)\}) \times \R^3
}
{ \longrightarrow} { \R^3 \setminus \{(0,0,0)\}
}
{ } {
}
}
{}{}{}
tentukan parameter-parameter
\mathl{(r,s,t)}{,} yang untuknya vektor $\left(  3   , \,   7  , \,   4                 \right)$ termasuk dalam serat
\mathl{L_{(r,s,t)}}{}.

}
{} {}

\inputaufgabe
{}
{

Deskripsikan
\definitionsverweis {bundel tangen}{}{}
dari
\definitionsverweis {sfera}{}{}
berdimensi $(n-1 )$

\mavergleichskettedisp
{\vergleichskette
{ S^{n-1}
}
{ =} { \left\{  \left(  x_1   , \,   \ldots  , \,   x_n          \right)   \mid   x_1^2 + \cdots + x_n^2 = 1         \right\}
}
{ \subset} { \R^n
}
{ } {
}
{ } {
}
}
{}{}{}
sebagai sebuah
\definitionsverweis {bundel kernel}{}{}
dalam arti
[[Topologischer Raum/Triviale Vektorbündel/Homomorphismus/Kernbündel/Bemerkung|Catatan 12.5]],
bandingkan pula dengan
[[Differenzierbare Hyperfläche/Tangentialbündel als Kernbündel/Beispiel|Contoh 12.7]].

}
{} {}

\inputaufgabe
{}
{

Misalkan $X$ sebuah
\definitionsverweis {ruang topologis}{}{.}
Tunjukkan bahwa sebuah
\definitionsverweis {homomorfisme}{}{}
antara bundel-bundel vektor
\zusatzklammer {trivial} {} {}
\maabbdisp {\varphi} {X \times \R^n } { X \times \R^m
} {}
sama dengan sebuah susunan berukuran
$m \times n$, yaitu sebuah
\definitionsverweis {matriks}{}{,}
yang entri-entrinya merupakan fungsi-fungsi kontinu dari $X$ ke $\R$.

}
{} {}

\inputaufgabe
{}
{

Misalkan

\mavergleichskette
{\vergleichskette
{U
}
{ \subseteq }{ \R^n
}
{  }{
}
{    }{
}
{   }{
}
}
{}{}{}
\definitionsverweis {terbuka}{}{}
dan
\maabb {f} { U } { \R^m
} {}
sebuah
\definitionsverweis {pemetaan terdiferensialkan kontinu}{}{.}
Tunjukkan bahwa
\definitionsverweis {diferensial total}{}{}
dalam bentuk
\maabbeledisp {} { U \times \R^n } { U \times \R^m
} { (x,v)} { (x, \left(D f \right)_{ x} (v) )
} {,}
memberikan sebuah homomorfisme dari bundel vektor
\mathl{U \times \R^n}{} ke bundel vektor
\mathl{U \times \R^m}{}.

}
{} {}

\inputaufgabe
{}
{

Perhatikan
\definitionsverweis {ruang topologis}{}{}

\mavergleichskettedisp
{\vergleichskette
{Y
}
{ \defeq} { \left\{ (s,t,u,v) \in \R^4  \mid  su+tv = 1        \right\}
}
{ } {
}
{ } {
}
{ } {
}
}
{}{}{}
dengan proyeksi
\maabbeledisp {p} { Y } { \R^2 \setminus \{ (0,0) \} = X
} {(s,t,u,v)} { (s,t)
} {.}
\aufzaehlungvier{Tunjukkan bahwa setiap serat $p$
\definitionsverweis {homeomorfik}{}{}
dengan sebuah garis riil.
}{Tunjukkan bahwa

\mavergleichskettedisp
{\vergleichskette
{ \varphi(s,t)
}
{ =} {(s,t,u(s,t),v(s,t))
}
{ =} { \left( s,t, { \frac{   s  }{ s^2+t^2 } } , { \frac{   t  }{ s^2+t^2 } }  \right)
}
{ } {
}
{ } {
}
}
{}{}{}
mendefinisikan sebuah pemetaan kontinu
\maabb {\varphi} { X} { Y
} {}
dengan

\mavergleichskettedisp
{\vergleichskette
{ p \circ \varphi
}
{ =} {
\operatorname{Id}_{  X  }
}
{ } {
}
{ } {
}
{ } {
}
}
{}{}{}.
}{Definisikan sebuah
\definitionsverweis {homeomorfisme}{}{}
antara
\mathkor {} {Y} {dan} {X \times \R} {.}
}{Tunjukkan bahwa tidak ada pemetaan polinomial
\maabb {\psi} { X } { Y
} {}
dengan

\mavergleichskettedisp
{\vergleichskette
{ p \circ \psi
}
{ =} {
\operatorname{Id}_{  X  }
}
{ } {
}
{ } {
}
{ } {
}
}
{}{}{}.
}

}
{} {}

\bild{ \begin{center}
\includegraphics[width=5.5cm]{ \bildeinlesungsvg {Inclusion-exclusion} {svg} }
\end{center}
\bildtext {Diagram Venn tiga himpunan A, B, dan C yang memperlihatkan daerah-daerah tumpang tindih untuk prinsip inklusi--eksklusi.} }

\bildlizenz { Inclusion-exclusion.svg } {Pembuat tidak diketahui} {} {} {domain publik} {}

Bukan hanya bundel vektor, tetapi secara lebih umum ruang-ruang topologis juga dapat dideskripsikan melalui data perekatan.

Yang dimaksud dengan sebuah
\definitionswort {data perekatan}{}
untuk
\definitionsverweis {ruang-ruang topologis}{}{}
adalah kumpulan data berikut.
\aufzaehlungvier{Sebuah keluarga ruang topologis

\mathbed {U_i} {}
 {i \in I} {}
 {} {} {} {,}
}{Untuk setiap pasangan
\mathl{(i,j)}{} sebuah
\definitionsverweis {himpunan bagian terbuka}{}{}

\mavergleichskette
{\vergleichskette
{U_{ij}
}
{ \subseteq }{ U_i
}
{  }{
}
{    }{
}
{   }{
}
}
{}{}{}
\zusatzklammer {dengan

\mavergleichskettek
{\vergleichskettek
{U_{ii}
}
{ = }{ U_i
}
{  }{
}
{    }{
}
{   }{
}
}
{}{}{}} {} {.}
}{Untuk setiap pasangan
\mathl{(i,j)}{} sebuah
\definitionsverweis {homeomorfisme}{}{}
\maabbdisp {\varphi_{ji}} {U_{ij}} { U_{ji}
} {}
\zusatzklammer {dengan

\mavergleichskettek
{\vergleichskettek
{ \varphi_{ii}
}
{ = }{
\operatorname{Id}_{ U_i  }
}
{  }{
}
{    }{
}
{   }{
}
}
{}{}{}} {} {.}
}{Untuk indeks-indeks

\mavergleichskette
{\vergleichskette
{ i,j,k
}
{ \in }{ I
}
{  }{
}
{    }{
}
{   }{
}
}
{}{}{}
terpenuhi
\definitionswort {syarat koklus}{}

\mavergleichskettedisp
{\vergleichskette
{ \varphi_{kj} \circ \varphi_{ji}
}
{ =} { \varphi_{ki}
}
{ } {
}
{ } {
}
{ } {
}
}
{}{}{}
sebagai pemetaan dari
\mathl{U_{ik} \cap U_{ij}}{} ke
\mathl{U_{k}}{}.
}

\inputaufgabegibtloesung
{}
{

Misalkan diberikan sebuah
\definitionsverweis {data perekatan}{}{}

\mathbed {U_i} {}
 {i \in I} {}
 {} {} {} {,}
untuk
\definitionsverweis {ruang-ruang topologis}{}{.}
Tunjukkan bahwa terdapat sebuah ruang topologis $X$ yang ditentukan secara tunggal, sebuah penutup terbuka

\mavergleichskette
{\vergleichskette
{ X
}
{ = }{ \bigcup_{i \in I} V_i
}
{  }{
}
{    }{
}
{   }{
}
}
{}{}{}
dan
\definitionsverweis {homeomorfisme-homeomorfisme}{}{}
\maabb {\psi_i} {U_i } { V_i
} {}
sedemikian sehingga

\mavergleichskettedisp
{\vergleichskette
{ \psi_i  \left(  U_{ij}  \right)
}
{ =} { V_i \cap V_j
}
{ } {
}
{ } {
}
{ } {
}
}
{}{}{}
dan

\mavergleichskettedisp
{\vergleichskette
{ \psi_i \vert_{U_{ij} }
}
{ =} { \psi_j \vert_{U_{ji} }  \circ   \varphi_{ji}
}
{ } {
}
{ } {
}
{ } {
}
}
{}{}{}.

}
{} {}

\inputaufgabegibtloesung
{}
{

Misalkan diberikan sebuah
\definitionsverweis {data perekatan}{}{}

\mathbed {E_i} {}
 {i \in I} {}
 {} {} {} {,}
di atas sebuah
\definitionsverweis {ruang topologis}{}{}

\mavergleichskettedisp
{\vergleichskette
{ X
}
{ =} { \bigcup_{i \in I} U_i
}
{ } {
}
{ } {
}
{ } {
}
}
{}{}{}.
Tunjukkan bahwa terdapat sebuah bundel vektor riil yang ditentukan secara tunggal
\maabb {} { E } { X
} {}
dan
\definitionsverweis {isomorfisme-isomorfisme}{}{}
\maabb {\psi_i} {E_i } { E \vert_{U_i }
} {}
sedemikian sehingga

\mavergleichskettedisp
{\vergleichskette
{ \psi_i \vert_{ E_i {{ |}}_{ U_{i} \cap U_j } }
}
{ =} { \psi_j \vert_{ E_j {{ |}}_{ U_{i} \cap U_j } }  \circ \varphi_{ji}
}
{ } {
}
{ } {
}
{ } {
}
}
{}{}{}.

}
{} {}

\inputaufgabe
{}
{

Misalkan
\maabb {s} { X } { V
} {}
sebuah
\definitionsverweis {penampang kontinu}{}{}
dari sebuah
\definitionsverweis {bundel vektor riil}{}{}
\maabb {p} { V } { X
} {}
di atas sebuah
\definitionsverweis {ruang topologis}{}{}
$X$. Tunjukkan bahwa
\definitionsverweis {citra}{}{}

\mavergleichskette
{\vergleichskette
{s(X)
}
{ \subseteq }{ V
}
{  }{
}
{    }{
}
{   }{
}
}
{}{}{}
merupakan sebuah
\definitionsverweis {himpunan bagian tertutup}{}{}
yang
\definitionsverweis {homeomorfik}{}{}
dengan $X$.

}
{} {}

\inputaufgabe
{}
{

Misalkan
\maabb {p} { V } { X
} {}
sebuah
\definitionsverweis {bundel vektor riil}{}{}
ber-
\definitionsverweis {rank}{}{}
$r$ di atas sebuah
\definitionsverweis {ruang topologis}{}{}
$X$, dan misalkan

\mavergleichskette
{\vergleichskette
{ U
}
{ \subseteq }{ X
}
{  }{
}
{    }{
}
{   }{
}
}
{}{}{}
sebuah
\definitionsverweis {himpunan bagian terbuka}{}{.}
Tunjukkan bahwa pada $U$ terdapat sebuah
\definitionsverweis {trivialisasi}{}{}
kontinu dari $V \vert_U$ jika dan hanya jika terdapat sebuah keluarga
\definitionsverweis {penampang-penampang basis}{}{}
kontinu
\maabbdisp {s_1 , \ldots , s_r} { U} {V  \vert_U
} {}.

}
{} {}

\inputaufgabe
{}
{

Misalkan sebuah
\definitionsverweis {bundel vektor}{}{}
\maabb {} { V } { X
} {}
ber-rank $m$ di atas sebuah
\definitionsverweis {ruang topologis}{}{}
$X$ diberikan oleh pemetaan-pemetaan matriks kontinu
\maabbdisp {\varphi_{ij}} { U_i \cap U_j } { \operatorname{GL}_{  m   } \! \left(  \R  \right)
} {}
untuk sebuah penutup terbuka

\mavergleichskette
{\vergleichskette
{X
}
{ = }{ \bigcup_{i \in I} U_i
}
{  }{
}
{    }{
}
{   }{
}
}
{}{}{}.
Tunjukkan bahwa sebuah
\definitionsverweis {penampang kontinu}{}{}
\maabbdisp {s} { X } { V
} {}
sama dengan sebuah keluarga pemetaan kontinu
\maabb {t_i} {U_i } { \R^m
} {,}
yang memenuhi syarat

\mavergleichskettedisp
{\vergleichskette
{t_j
}
{ =} { \varphi_{ij} t_i
}
{ } {
}
{ } {
}
{ } {
}
}
{}{}{}
untuk semua $i,j$.

}
{} {}

\inputaufgabe
{}
{

Misalkan

\mavergleichskette
{\vergleichskette
{X
}
{ = }{ \bigcup_{i \in I} U_i
}
{  }{
}
{    }{
}
{   }{
}
}
{}{}{}
sebuah
\definitionsverweis {penutup terbuka}{}{}
dari sebuah
\definitionsverweis {ruang topologis}{}{}
$X$. Misalkan
\maabbdisp {\varphi_{ij}} {U_i \cap U_j } { \operatorname{GL}_{  r   } \! \left(  \R  \right)
} {}
dan
\maabbdisp {\psi_{ij}} {U_i \cap U_j } { \operatorname{GL}_{  r   } \! \left(  \R  \right)
} {}
merupakan
\definitionsverweis {deskripsi-deskripsi matriks}{}{}
yang secara berturut-turut menghasilkan bundel vektor
\mathkor {} {E} {dan} {F} {.}
Tunjukkan bahwa bundel-bundel ini
\definitionsverweis {isomorfik}{}{}
jika dan hanya jika terdapat pemetaan-pemetaan kontinu
\maabbdisp {\alpha_i} { U_i } { \operatorname{GL}_{  r   } \! \left(  \R  \right)
} {}
sedemikian sehingga
\zusatzklammer {untuk pembatasan-pembatasan yang bermakna} {} {}

\mavergleichskettedisp
{\vergleichskette
{ \alpha_i \circ  \varphi_{ij}  \circ \alpha_j^{-1}
}
{ =} { \psi_{ij}
}
{ } {
}
{ } {
}
{ } {
}
}
{}{}{}
untuk semua $i,j$.

}
{} {}

\inputaufgabe
{}
{

Tunjukkan bahwa komplemen penampang nol dalam sebuah
\definitionsverweis {bundel garis}{}{}
trivial tidak
\definitionsverweis {terhubung}{}{.}

}
{} {}

\inputaufgabe
{}
{

Ambillah sebuah pita persegi panjang yang tipis, puntir pita itu satu putaran penuh terhadap sumbu panjangnya, lalu rekatkan kedua sisi pendeknya. Sekarang potong pita itu memanjang tepat di tengah dengan gunting. Apakah objek yang dihasilkan
\definitionsverweis {terhubung}{}{?}
Bagaimanakah bentuknya jika dibuat $n$ setengah putaran?

}
{} {}

Misalkan
\maabb {} { V } { X
} {}
sebuah
\definitionsverweis {bundel vektor riil}{}{}
di atas sebuah
\definitionsverweis {ruang topologis}{}{.}
Sebuah bundel vektor riil
\maabb {} { U } { X
} {}
disebut
\definitionswort {subbundel}{}
dari $V$ jika diberikan sebuah
\definitionsverweis {homomorfisme bundel vektor}{}{}
\definitionsverweis {injektif}{}{}
\maabb {} { U } { V
} {}.

\inputaufgabe
{}
{

Misalkan $X$ sebuah
\definitionsverweis {ruang topologis}{}{}
dan
\maabb {\varphi} { V } { W
} {}
sebuah
\definitionsverweis {homomorfisme}{}{}
antara
\definitionsverweis {bundel-bundel vektor riil}{}{}
di atas $X$. Misalkan

\mavergleichskettedisp
{\vergleichskette
{ Y
}
{ =} { \left\{ P \in X  \mid  \varphi_P \text{ surjektif }         \right\}
}
{ } {
}
{ } {
}
{ } {
}
}
{}{}{}.
Tunjukkan bahwa
\zusatzklammer {keluarga} {} {}
\definitionsverweis {subruang-subruang vektor}{}{}

\mavergleichskette
{\vergleichskette
{ \operatorname{kern} \varphi_P
}
{ \subseteq }{ V_P
}
{  }{
}
{    }{
}
{   }{
}
}
{}{}{}
membentuk sebuah
\definitionsverweis {subbundel vektor}{}{}
dari $V\vert_Y$.

}
{} {}

\inputaufgabe
{}
{

Misalkan $X$ sebuah
\definitionsverweis {ruang topologis}{}{}
dan

\mavergleichskette
{\vergleichskette
{ U
}
{ \subseteq }{ V
}
{  }{
}
{    }{
}
{   }{
}
}
{}{}{}
sebuah
\definitionsverweis {subbundel vektor}{}{}
dari
\definitionsverweis {bundel vektor riil}{}{}
\maabb {} { V } { X
} {.}
Tunjukkan bahwa
\zusatzklammer {keluarga} {} {}
\definitionsverweis {ruang-ruang hasil bagi}{}{}

\mathl{V_P/U_P}{} membentuk sebuah bundel vektor di atas $X$ dan bahwa terdapat sebuah
\definitionsverweis {homomorfisme bundel}{}{}
\definitionsverweis {surjektif}{}{}
\maabb {} { V } { V/U
} {}.

}
{} {}

Kita perkenalkan dua istilah lagi.

Misalkan $K$ sebuah
\definitionsverweis {lapangan}{}{}
dan
\mathl{U,V,W}{}
\definitionsverweis {ruang-ruang vektor}{}{}
di atas $K$. Sebuah diagram berbentuk

\mathdisp {0 \, \stackrel{  }{\longrightarrow} \, U  \, \stackrel{  }{ \longrightarrow}   \, V  \, \stackrel{  }{\longrightarrow} \, W \,\stackrel{  } {\longrightarrow}  \, 0} {  }

disebut sebuah \definitionswort {barisan eksak pendek}{} ruang-ruang vektor-$K$ jika $U$ merupakan sebuah
\definitionsverweis {subruang vektor}{}{}
dari $V$ dan jika $W$
\definitionsverweis {isomorfik}{}{}
dengan
\definitionsverweis {ruang hasil bagi}{}{}
$V/U$.

Misalkan $X$ sebuah
\definitionsverweis {ruang topologis}{}{,}
misalkan
\maabbdisp {} {U,V,W} {X
} {}
\definitionsverweis {bundel-bundel vektor}{}{}
di atas $X$, serta misalkan
\maabb {\varphi} { U } { V
} {}
dan
\maabb {\psi} { V } { W
} {}
\definitionsverweis {homomorfisme-homomorfisme}{}{.}
Kita katakan bahwa terdapat sebuah
\definitionswort {barisan eksak pendek}{}

\mathdisp {0 \, \stackrel{  }{\longrightarrow} \,  U   \, \stackrel{ \varphi }{ \longrightarrow}   \,  V   \, \stackrel{ \psi }{\longrightarrow} \,  W  \,\stackrel{  } {\longrightarrow}  \, 0} {  }

jika untuk setiap titik

\mavergleichskette
{\vergleichskette
{ P
}
{ \in }{ X
}
{  }{
}
{    }{
}
{   }{
}
}
{}{}{}
terdapat sebuah
\definitionsverweis {barisan eksak pendek}{}{}

\mathdisp {0 \, \stackrel{  }{\longrightarrow} \, U_P   \, \stackrel{ \varphi_P }{ \longrightarrow}   \,  V_P   \, \stackrel{ \psi_P }{\longrightarrow} \,  W_P  \,\stackrel{  } {\longrightarrow}  \, 0} {  }

pada serat-seratnya.

\inputaufgabe
{}
{

Misalkan $X$ sebuah
\definitionsverweis {ruang topologis}{}{}
dan

\mavergleichskette
{\vergleichskette
{ U
}
{ \subseteq }{ V
}
{  }{
}
{    }{
}
{   }{
}
}
{}{}{}
sebuah
\definitionsverweis {subbundel vektor}{}{}
dari
\definitionsverweis {bundel vektor riil}{}{}
\maabb {} { V } { X
} {.}
Tunjukkan bahwa terdapat sebuah
\definitionsverweis {barisan eksak pendek}{}{}

\mathdisp {0 \, \stackrel{  }{\longrightarrow} \,  U   \, \stackrel{  }{ \longrightarrow}   \,  V   \, \stackrel{  }{\longrightarrow} \,  V/U \,\stackrel{  } {\longrightarrow}  \, 0} {  }.

}
{} {}

\inputaufgabe
{}
{

Misalkan
\mathkor {} {X} {dan} {Y} {}
\definitionsverweis {ruang-ruang topologis}{}{}
dan
\maabb {\varphi} { Y } { X
} {}
sebuah
\definitionsverweis {pemetaan kontinu}{}{.}
Misalkan
\maabb {p} { V } { X
} {}
sebuah
\definitionsverweis {bundel vektor}{}{}
di atas $X$. Tunjukkan bahwa
\mathl{Y \times_X V}{}
\zusatzklammer {lihat
[[Topologische Räume/Relatives Produkt/Aufgabe|Soal 11.16]]} {} {}
merupakan sebuah bundel vektor di atas $Y$.

}
{} {}

\inputaufgabe
{}
{

Misalkan
\mathkor {} {X} {dan} {Y} {}
\definitionsverweis {ruang-ruang topologis}{}{}
dan
\maabb {\varphi} {Y} {X
} {}
sebuah
\definitionsverweis {pemetaan kontinu}{}{.}
Misalkan
\maabb {p} {V} {X
} {}
sebuah
\definitionsverweis {bundel vektor}{}{}
di atas $X$. Deskripsikan
\definitionsverweis {tarik balik}{}{}
sebuah bundel vektor dengan bantuan
\definitionsverweis {data perekatan}{}{.}

}
{} {}

\inputaufgabe
{}
{

Misalkan $X$ sebuah
\definitionsverweis {ruang topologis}{}{}
dan
\maabb {p} { V } { X
} {}
serta
\maabb {q} { W } { X
} {}
\definitionsverweis {bundel-bundel vektor}{}{}
di atas $X$. Tunjukkan bahwa
\mathl{V \times_X W}{}
\zusatzklammer {lihat
[[Topologische Räume/Relatives Produkt/Aufgabe|Soal 11.16]]} {} {}
merupakan sebuah bundel vektor di atas $X$ yang sama dengan
\definitionsverweis {jumlah langsung bundel-bundel vektor}{}{}
di atas $X$.

}
{} {}

\inputaufgabe
{}
{

Misalkan
\maabb {\varphi} { X } { Y
} {}
sebuah
\definitionsverweis {pemetaan diferensiabel}{}{}
antara
\definitionsverweis {manifold-manifold diferensiabel}{}{}
\mathkor {} {X} {dan} {Y} {.}
Tunjukkan bahwa pemetaan ini menghasilkan
\definitionsverweis {homomorfisme bundel vektor}{}{}
alami
\maabbdisp {T\varphi} {TX} { \varphi^*TY
} {}
antara
\definitionsverweis {bundel-bundel vektor}{}{}
di atas $X$.

}
{} {}

  \zwischenueberschrift{Soal untuk dikumpulkan}

\inputaufgabe
{4}
{

Misalkan $X$ sebuah
\definitionsverweis {ruang topologis}{}{}
dan misalkan

\mavergleichskette
{\vergleichskette
{ f_1 , \ldots , f_n
}
{ \in }{ C^0(X,\R)
}
{  }{
}
{    }{
}
{   }{
}
}
{}{}{}
merupakan fungsi-fungsi bernilai riil
\definitionsverweis {kontinu}{}{}
pada $X$. Misalkan

\mavergleichskette
{\vergleichskette
{ U
}
{ = }{ \bigcup_{i = 1}^n D(f_i)
}
{ \subseteq }{ X
}
{    }{
}
{   }{
}
}
{}{}{}
dan

\mavergleichskettedisp
{\vergleichskette
{ S
}
{ =} { \left\{ (P; t_1 , \ldots , t_n)  \mid  P \in U , \,  \sum_{i = 1}^n f_i(P)t_i = 0       \right\}
}
{ \subseteq} { U \times \R^n
}
{ } {
}
{ } {
}
}
{}{}{,}
bandingkan dengan
[[Topologischer Raum/Triviale Vektorbündel/Homomorphismus/Kernbündel/Bemerkung|Catatan 12.5]].
Berikan secara eksplisit trivialisasi-trivialisasi untuk $S$ di atas $D(f_i)$ dan tunjukkan bahwa $S$ merupakan sebuah bundel vektor di atas $U$ ber-rank $n-1$.

}
{} {}

\inputaufgabe
{4}
{

Tunjukkan bahwa sebuah
\definitionsverweis {bundel garis riil}{}{}
\maabb {} { L } { X
} {}
di atas sebuah
\definitionsverweis {ruang topologis}{}{}
$X$ adalah trivial jika dan hanya jika bundel itu memiliki sebuah
\definitionsverweis {penampang}{}{}
kontinu yang tidak pernah nol.

}
{} {}

\inputaufgabe
{3}
{

Tunjukkan bahwa sebuah data perekatan untuk sebuah bundel garis terhadap sebuah penutup terbuka

\mavergleichskette
{\vergleichskette
{X
}
{ = }{ \bigcup_{i \in I} U_i
}
{  }{
}
{    }{
}
{   }{
}
}
{}{}{}
sama dengan sebuah keluarga fungsi kontinu yang tidak pernah nol
\maabb {f_{ij}} {U_i \cap U_j } { \R
} {,}
yang pada
\mathl{U_i \cap U_j \cap U_k}{} masing-masing memenuhi syarat

\mavergleichskette
{\vergleichskette
{ f_{ki}
}
{ = }{ f_{kj} \cdot f_{ji}
}
{  }{
}
{    }{
}
{   }{
}
}
{}{}{.}

}
{} {}

\inputaufgabe
{4}
{

Tunjukkan bahwa
\definitionsverweis {pita Möbius}{}{}
merupakan sebuah
\definitionsverweis {manifold}{}{}
berdimensi dua.

}
{} {}

 [[Kategorie:Latexseite]]
